\documentclass[10pt]{article}
\usepackage[T1]{fontenc}
\usepackage[utf8]{inputenc}
\usepackage[french]{babel}
\usepackage{amsmath,amssymb,mathrsfs,mathpazo}
\usepackage{geometry}
\geometry{a4paper,margin=2cm}
\usepackage{graphicx,array,tabularx,tabularray,colortbl,multirow,multicol,booktabs,enumitem,pifont,ulem,fancybox,diagbox,cancel}
\usepackage{pstricks,pstricks-add,pst-plot,pst-tree,pst-node,pst-text,pst-func,pst-math,pst-bspline,pst-3dplot,pst-eucl,pst-solides3d}
\let\pscalculate\undefined
\usepackage{pst-bezier}
\usepackage{tikz,pgfplots,tkz-tab}
\usetikzlibrary{arrows,arrows.meta,patterns,calc,positioning,trees,decorations.markings,shapes.geometric,shapes.misc,shapes.arrows,intersections,decorations.pathmorphing,decorations.pathreplacing,backgrounds,shadows,babel}
\pgfplotsset{compat=1.18}
\usepgfplotslibrary{fillbetween}
\usepackage[np]{numprint}
\usepackage{listings,makecell,pas-tableur,twoopt}
\let\barre\undefined
\usepackage[pstricks]{tablvar}
\definecolor{darkgreen}{RGB}{0,100,0}
\providecommand{\red}{\color{red}}
\providecommand{\blue}{\color{blue}}
\providecommand{\green}{\color{green}}
\providecommand{\black}{\color{black}}
\providecommand{\white}{\color{white}}
\providecommand{\gray}{\color{gray}}
\providecommand{\cyan}{\color{cyan}}
\providecommand{\magenta}{\color{magenta}}
\providecommand{\yellow}{\color{yellow}}
\providecommand{\violet}{\color{violet}}
\providecommand{\orange}{\color{orange}}
\providecommand{\euro}{\texteuro}
\providecommand{\EUR}{\texteuro}
\providecommand{\ieme}{\textsuperscript{e}}
\providecommand{\degres}{\ensuremath{^\circ}}
\providecommand{\degree}{\ensuremath{^\circ}}
\providecommand{\vv}[1]{\overrightarrow{#1}}
\providecommand{\Coord}[1]{#1}
\providecommand{\decosix}{}
\providecommand{\starredbullet}{\ensuremath{\bullet}}
\pagestyle{empty}
\setlength{\parindent}{0pt}
\usepackage{adjustbox}\DeclareUnicodeCharacter{2236}{:}\DeclareUnicodeCharacter{2190}{\ensuremath{\leftarrow}}
\providecommand{\R}{\mathbb{R}}
\providecommand{\N}{\mathbb{N}}
\providecommand{\Z}{\mathbb{Z}}
\providecommand{\Q}{\mathbb{Q}}
\providecommand{\D}{\mathbb{D}}
\providecommand{\C}{\mathbb{C}}
\providecommand{\e}{{\mathrm{e}}}
\providecommand{\vect}[1]{\overrightarrow{#1}}
\providecommand{\vectt}[1]{\overrightarrow{\mathrm{#1}}}
\providecommand{\vvt}[1]{\overrightarrow{\mathrm{#1}}}
\providecommand{\np}[2][]{{#2~#1}}
\providecommand{\vv}[1]{{\overrightarrow{#1}}}
\renewcommand{\d}{{\mathrm{d}}}
\providecommand{\degres}{{^\circ}}
\providecommand{\degre}{{^\circ}}
\providecommand{\blue}{\color{blue}}
\providecommand{\red}{\color{red}}
\providecommand{\enskip}{\quad}
\providecommand{\llbracket}{\lbrack\!\lbrack}
\providecommand{\rrbracket}{\rbrack\!\rbrack}
\providecommand{\euro}{{\text{€}}}
\providecommand{\no}{{\text{n°}}}
\providecommand{\interval}[2]{{[#1;#2]}}
\providecommand{\intervFO}[2]{{[#1;#2[}}
\providecommand{\dotfill}{\ldots}

\providecommand{\ds}{\displaystyle}
\providecommand{\textding}[1]{#1}
\providecommand{\Oij}{\(\mathrm{O};\vec{\imath},\vec{\jmath}\)}
\providecommand{\Oijk}{\(\mathrm{O};\vec{\imath},\vec{\jmath},\vec{k}\)}
\providecommand{\Ouv}{\(\mathrm{O};\vec{u},\vec{v}\)}
\providecommand{\intervFF}[2]{\ensuremath{[#1;#2]}}
\renewcommand{\euro}{\ifmmode\text{\texteuro}\else\texteuro\fi}
\DeclareUnicodeCharacter{00B0}{\ifmmode^{\circ}\else\textdegree\fi}
\DeclareUnicodeCharacter{20AC}{\ifmmode\text{\texteuro}\else\texteuro\fi}
% Énoncés : Éducation nationale. Transcription et corrigés : APMEP, Denis Vergès et contributeurs. Découpage et classement : CoopMaths, Sébastien Lozano, Éric Elter et contributeurs. Adaptation de la présentation : Maths Explorer.
% CC BY-SA : https://coopmaths.fr/www/about/
% https://coopmaths.fr/alea/static/eam/2026/tex/eam_technologique_2026_06_antilles_1.tex
\begin{document}
\noindent\textit{APMEP / CoopMaths. CC BY-SA. Présentation adaptée par Maths Explorer.}\par\medskip
\begin{minipage}[t]{0.6\linewidth}
	Au $1{ }^{\text {er }}$ janvier 2025, une famille installe des éoliennes domestiques qui produiront une quantité d'énergie de \np[kWh]{3000} en 2025.

	On constate que chaque année, la quantité d'énergie produite par cette installation baisse de 5\,\%.

	Pour tout entier $n \geqslant  0$, on note $u_{n}$ la quantité d'énergie (en kWh) produite par l'installation durant l'année $(2025+n)$.

	On a donc $u_{0}=3000$.

	\begin{enumerate}
		\item \begin{enumerate}
			\item Quelle sera la quantité d'énergie produite par l'installation durant l'année 2026 ?

			Détailler les calculs.

			\item Justifier que, pour tout entier naturel $n$,
			\[u_{n+1}=0,95 \times u_{n}.\]

			\item Quelle est la nature de la suite $(u_{n})$ ?

			Préciser sa raison.
		\end{enumerate}
	\end{enumerate}

	À l'aide d'un tableur, on a obtenu les valeurs successives des termes de la suite $(u_{n})$.

	On donne ci-contre un extrait de la feuille de calculs.

	\begin{enumerate}[resume]
		\item Vous répondrez aux questions suivantes à l'aide du tableau ci-contre.

		\begin{enumerate}
			\item Déterminer à partir de quelle année la quantité d'énergie produite sera inférieure à \np[kWh]{1000}.

			\item Après 25 ans de fonctionnement, le fabricant affirme que la quantité d'énergie produite s'élèvera à plus de $30\,\%$ de la quantité produite initialement.

			Cette affirmation est-elle vraie ou fausse ?

			Justifier.
		\end{enumerate}
	\end{enumerate}
	\end{minipage}
	\hfill
	\begin{tblr}[T]{hlines, vlines, rowsep=0.5pt, columns={c,font=\sffamily},column{1}={wd=7mm,colsep=2pt},
		row{1}={rowsep=0pt,bg=gray9,rowsep =0pt},
		column{1}={colsep=0pt,7mm,bg=gray9},
		cell{1}{1}={gray9,halign=r,valign=b,cmd={\tikz{\fill[gray!50] (0,0)--(0.4,0.4)--(0.4,0);\clip (0,0) rectangle (0.5,0.45);}}}}
& A           & B                                   \\
		1  & Valeur de n & Valeur de u\textsubscript{n} \\
		2  & 0           & 3000,00                             \\
		3  & 1           &                                     \\
		4  & 2           & 2707,50                             \\
		5  & 3           & 2572,13                             \\
		6  & 4           & 2443,52                             \\
		7  & 5           & 2321,34                             \\
		8  & 6           & 2205,28                             \\
		9  & 7           & 2095,01                             \\
		10 & 8           & 1990,26                             \\
		11 & 9           & 1890,75                             \\
		12 & 10          & 1796,21                             \\
		13 & 11          & 1706,40                             \\
		14 & 12          & 1621,08                             \\
		15 & 13          & 1540,03                             \\
		16 & 14          & 1463,02                             \\
		17 & 15          & 1389,87                             \\
		18 & 16          & 1320,38                             \\
		19 & 17          & 1254,36                             \\
		20 & 18          & 1191,64                             \\
		21 & 19          & 1132,06                             \\
		22 & 20          & 1075,46                             \\
		23 & 21          & 1021,68                             \\
		24 & 22          & 970,60                              \\
		25 & 23          & 922,07                              \\
		26 & 24          & 875,97                              \\
		27 & 25          & 832,17                              \\
		28 & 26          & 790,56                              \\
		29 & 27          & 751,03                              \\
		30 & 28          & 713,48                              \\
		31 & 29          & 677,81                              \\
		32 & 30          & 643,92                              \\
		33 & 31          & 611,72                              \\
		34 & 32          & 581,13
	\end{tblr}

	\begin{enumerate}
		\setcounter{enumi}{2}
		\item On a représenté ci-dessous deux suites.

		Indiquer le graphique correspondant à la représentation de la suite $\left(u_n\right) $.

		Justifier votre réponse.

		\begin{tikzpicture}[x = 4mm,y = 0.03mm,> =stealth]
			\draw[->] (-8mm,0)--(32,0) ;
			\foreach \x in {10,20,30}
			{\draw (\x,2pt)--(\x,-2pt) node[below]{\np{\x}} ;}
			\draw[->] (0,-8mm)--(0,3200);
			\foreach \y in {500,1000,...,3000}
			{\draw (2pt,\y)--(-2pt,\y) node[left]{\np{\y}} ;}
			\node[below left] at (-2pt,-2pt) {0};
			\foreach \n in {0,...,29}
			{\node[red] at (\n,{3000-75*\n}) {\textbf{+}};}
			\node[left] at(30,2400) {Graphique 1};
		\end{tikzpicture}

		\begin{tikzpicture}[x = 4mm,y = 0.03mm,> =stealth]
			\draw[->] (-8mm,0)--(32,0) ;
			\foreach \x in {10,20,30}
			{\draw (\x,2pt)--(\x,-2pt) node[below]{\np{\x}} ;}
			\draw[->] (0,-8mm)--(0,3200);
			\foreach \y in {500,1000,...,3000}
			{\draw (2pt,\y)--(-2pt,\y) node[left]{\np{\y}} ;}
			\node[below left] at (-2pt,-2pt) {0};
			\foreach \n in {0,...,29}
			{\node[red] at (\n,{3000*0.95^\n}) {\textbf{+}};}
			\node[left] at(30,2400) {Graphique 2};
		\end{tikzpicture}
	\end{enumerate}

\end{document}
